\hypertarget{o10strength:a-strengthened-opposite-chirality-gap-and-o10-through-euler-density-394}{%
\section{A strengthened opposite-chirality gap and O10 through Euler
density
39/4}\label{o10strength:a-strengthened-opposite-chirality-gap-and-o10-through-euler-density-394}}

2026-09-29. This proof strengthens the earlier critical-density result
by combining the bounded resolvent section with its full actual
first-torsion projection, a Hopf-form identity, and a sharp
three-dimensional spectral inequality. The earlier \(3/8\) theorem
remains a shorter-dependency special case. The strengthened theorem uses
the full strong-current conclusion \(N_\tau\to0\), but does not use the
subsequent first-prolongation estimate.

\hypertarget{o10strength:statement}{%
\subsection{1. Statement}\label{o10strength:statement}}

Retain O10's actual conformal monopole supply, \(p\ge A>0\), and
\(a^2\ge T_+\). There is an absolute \(\eta_*>0\) such that every actual
scalar-minus-four input and every \(\beta\ge4\) satisfy

\[
 \boxed{\kappa_\beta+\beta^{-2}+
 \left(\fint|W^-|^2-\frac{50}{3}\right)_+\ge\eta_*},
 \qquad \kappa_\beta=P/V+D_\beta/V.
\tag{0}
\]

This primary theorem has no Euler-density hypothesis. Consequently there
are absolute \(\delta_*>0\) and \(c_*>0\) such that

\[
 \frac{\chi(M)}A\le\frac{39}{4}+\delta_*
 \quad\Longrightarrow\quad
 \boxed{\fint|\operatorname{Ric}+g|^2+\frac{D_\beta}V
 \ge c_*\frac{(\beta-2)^2}{\beta^2}\quad(\beta\ge2)}.
\tag{1}
\]

The second statement holds for all original admissible metrics. No
pointwise curvature, diameter, Sobolev, noncollapse, volume, or
injectivity-radius bound is added. The constants are existential
absolute constants; numerical values for their positive margins are not
asserted.

If (0) failed, its three nonnegative summands would give a sequence
satisfying

\[
 \beta_j\to\infty,\qquad\kappa_j\to0,\qquad
 b_j:=\fint|W_j^-|^2\le50/3+o(1).
\tag{2a}
\]

We exclude precisely this sequence. Its opposite-Weyl energy is bounded
from the start; no Euler-density assumption is needed in the sequential
proof.

\hypertarget{o10strength:inputs-from-the-actual-endpoint-system}{%
\subsection{2. Inputs from the actual endpoint
system}\label{o10strength:inputs-from-the-actual-endpoint-system}}

All integral norms in this proof use \(dV/V\). Work with one actual
solution

\[
 s=-4,\qquad D_A\Phi=0,\qquad F_A^+=2i\sigma(\Phi),
 \qquad \rho=|\Phi|^2\le2.
\]

Write \(\tau=t_\Phi\in\Gamma(E)\), where \(E=S^-\otimes L\), and put

\[
 m=|\tau|^2,\quad \rho_*=4/3,\quad m_*=128/81,
 \quad S=m_*+2\rho^2|\nabla_A\Phi|^2,
 \quad Q=-iq_-(\tau),
\]

\[
 f=-iF_A^-,\qquad K=-\frac98\frac{\rho^3Q}{S}.
\]

The original manuscript proves \(|K|\le C\), \(f-K\to0\) in normalized
\(L^2\), and the scalar/curvature/current budgets. The strong-current
upgrade proved on 2026-09-29 additionally gives

\[
 \rho\to\rho_*\text{ in }L^2,\quad
 m\to m_*\text{ in }L^2,\quad
 D^-_{3A}\tau\to0\text{ in }L^2,\quad
 \|\nabla_{3A}\tau\|_2\le C.
\tag{2}
\]

The slow-gradient theorem for the original positive field first gives
the short chain

\[
 |\nabla_A\Phi|^4\text{ uniformly integrable}
 \Longrightarrow D\tau\to0\text{ in }L^2,\quad m\to m_*\text{ in }L^2.
\]

The strengthened argument additionally uses the complete previously
established strong-current chain, ending in

\[
 N_\tau=(\nabla\tau)^{1,0}_{J_+}\longrightarrow0
 \quad\text{in normalized }L^2.
\tag{2b}
\]

It is enough to use this projection on sets where rho is bounded away
from zero. The omitted low-rho regions have relative volume tending to
zero and will be controlled by the fourth-power gradient bound of the
resolvent section. The strong-current chain uses the positive-spinor
curvature blocks and the actual curvature fit; it does not assume
smallness of the opposite-Weyl block. The earlier \(3/8\) proof did not
need (2b); the present stronger statement does.

Let \(D=D^-_{3A}\). Its Lichnerowicz potential is

\[
 D^*D=T+\mathcal V,\quad
 T=\nabla_{3A}^*\nabla_{3A},\quad
 \mathcal V=-I+\frac32c(F_A^-).
\tag{3}
\]

The connection curvature of \(E\) is

\[
 F^E=F^{S^-}_{\rm LC}+\frac32F_A I.
\]

Initially (2a) gives only a uniform normalized L2 bound on the full
curvature of E. The scalar block and determinant curvature model are
bounded, while the opposite-Weyl block is L2 bounded. After proving the
sharp spectral constraint in Section 6, we shall improve this to

\[
 F^E=F_0+\mathcal E_F,\qquad |F_0|\le C,
 \qquad\|\mathcal E_F\|_{2,\rm avg}\to0.
\tag{4}
\]

It would be circular to assume (4) before that spectral step. The
resolvent construction in Sections 3--5 does not require it.

\hypertarget{o10strength:a-bounded-truncation-with-small-dirac-residual}{%
\subsection{3. A bounded truncation with small Dirac
residual}\label{o10strength:a-bounded-truncation-with-small-dirac-residual}}

Fix \(L>m_*\). Choose a smooth nondecreasing \(\varphi\), zero for
\(m\le L\), one for \(m\ge2L\), with \(m\varphi'\) bounded. Test (3)
against \(\varphi(m)\tau\). The exact identity is

\[
 \int\varphi|\nabla\tau|^2+
 \frac12\int\varphi'|dm|^2
 =\int\varphi|D\tau|^2+
 \int\operatorname{Re}\langle D\tau,c(d\varphi)\tau\rangle
 -\int\varphi\langle\mathcal V\tau,\tau\rangle.
\tag{5}
\]

The mixed term is bounded by

\[
 \frac14\int\varphi'|dm|^2+
 C\int m\varphi'|D\tau|^2.
\]

The potential term tends to zero after division by \(V\). For its
bounded-model part, \(m\to m_*\) in \(L^2\) and \(L>m_*\) imply
\(\fint_{m>L}m\to0\). Its remaining part is bounded by
\(C\|f-K\|_{2,\rm avg}\|m\|_{2,\rm avg}\to0\). Thus (2) and (5) give

\[
 \fint_{m\ge2L}|\nabla\tau|^2\longrightarrow0.
\tag{6}
\]

Choose a smooth radial cap \(b(m)\) equal to one on \([0,2L]\), equal to
a fixed constant times \(m^{-1/2}\) above \(3L\), and with
\(|b|+m|b'|\le C_L\). Set

\[
 u=b(m)\tau.
\]

This is a globally smooth section, uniformly bounded in norm. Equations
(2), (6) and the product rule give

\[
 \|u-\tau\|_{2,\rm avg}+\|\nabla(u-\tau)\|_{2,\rm avg}\to0,
 \qquad \|Du\|_{2,\rm avg}\to0,
 \qquad |u|^2\to m_*\text{ in }L^2_{\rm avg}.
\tag{7}
\]

The stronger last assertion also follows directly from boundedness of
\(u\), (2), and convergence in measure.

\hypertarget{o10strength:the-potential-acts-on-this-bounded-section-as--2}{%
\subsection{\texorpdfstring{4. The potential acts on this bounded
section as
\(-2\)}{4. The potential acts on this bounded section as -2}}\label{o10strength:the-potential-acts-on-this-bounded-section-as--2}}

The cap \(u\) is collinear with \(\tau\). The exact Hopf identity
\(c(q_-(\tau))\tau=(m/2)\tau\) implies

\[
 \left[-I+\frac32c(iK)\right]u
 =-\left(1+\frac{27}{32}\frac{\rho^3m}{S}\right)u.
\tag{8}
\]

The displayed coefficient is uniformly bounded, because \(m\le S\) and
\(\rho\le2\). By the actual calibration and (2), it tends in measure to

\[
 1+\frac{27}{32}\frac{\rho_*^3}{2}=2.
\]

Boundedness of \(u\), together with \(f-K\to0\) in \(L^2\), therefore
gives

\[
 \boxed{\|(\mathcal V+2I)u\|_{2,\rm avg}\to0.}\tag{9}
\]

Equations (3), (7), and (9) say that \(Tu-2u\) tends to zero in the
actual normalized \(H^{-1}\) norm. This negative-norm statement is used
only in the following coercive resolvent equation.

\hypertarget{o10strength:rough-laplacian-resolvent-bounded-forcing-and-strong-h1-approximation}{%
\subsection{\texorpdfstring{5. Rough-Laplacian resolvent: bounded
forcing and strong \(H^1\)
approximation}{5. Rough-Laplacian resolvent: bounded forcing and strong H\^{}1 approximation}}\label{o10strength:rough-laplacian-resolvent-bounded-forcing-and-strong-h1-approximation}}

Define, on each actual closed manifold and bundle,

\[
 v=3(I+T)^{-1}u.
\tag{10}
\]

The positive operator \(I+T\) is invertible on a closed manifold. The
connection-Laplacian maximum principle (or its scalar Kato inequality)
gives

\[
 \|v\|_\infty\le3\|u\|_\infty.
\tag{11}
\]

In particular \(v\) is a smooth bounded section and

\[
 Tv=3u-v=:h,\qquad \|h\|_\infty\le C.
\tag{12}
\]

Put \(z=v-u\). From (3),

\[
 (I+T)z=2u-Tu=-D^*Du+(\mathcal V+2I)u.
\]

Pair with \(z\). Since pointwise \(|Dz|\le2|\nabla z|\),

\[
 \|z\|_{2,\rm avg}^2+\|\nabla z\|_{2,\rm avg}^2
 \le2\|Du\|_{2,\rm avg}\|\nabla z\|_{2,\rm avg}
 +\|(\mathcal V+2I)u\|_{2,\rm avg}\|z\|_{2,\rm avg}.
\]

Thus

\[
 \boxed{\|v-u\|_{H^1_{\rm avg}}\to0,\qquad
               |v|^2\to m_*\text{ in }L^2_{\rm avg}.}\tag{13}
\]

The same equations give the stronger residuals

\[
 \|Dv\|_{2,\rm avg}\to0,\qquad
 \|Tv-2v\|_{2,\rm avg}=3\|u-v\|_{2,\rm avg}\to0.\tag{13a}
\]

The first follows from \(v-u\to0\) in \(H^1\) and \(Du\to0\); the second
is the exact resolvent equation. These are valid even at arbitrary fixed
Euler density. The rough-Laplacian residual is genuinely \(L^2\), rather
than only \(H^{-1}\).

No spectral gap for \(T\) was used; the identity term in its resolvent
supplies a universal coercivity constant. No derivative of \(Du\), of
Ricci, or of bundle curvature was estimated.

\hypertarget{o10strength:bounded-rough-laplacian-suffices-for-critical-gradient-control}{%
\subsection{6. Bounded rough Laplacian suffices for critical gradient
control}\label{o10strength:bounded-rough-laplacian-suffices-for-critical-gradient-control}}

We record the exact extension of the slow-truncation lemma used here.
Let \(v\) be a smooth bounded section with bounded
\(h=\nabla^*\nabla v\). Suppose Ricci and bundle curvature are each a
uniformly bounded tensor plus an error whose normalized \(L^2\) norm
tends to zero. Then, for some fixed \(\epsilon>0\),

\[
 \fint\left(\frac{|\nabla v|^2}{R^2+|\nabla v|^2}\right)^\epsilon
       \bigl(|\nabla^2v|^2+|\nabla v|^4\bigr)
 \le CR^{-2\epsilon}+o(1),\quad R\ge1.
\tag{14}
\]

In particular \(|\nabla v|^4\) is uniformly integrable relative to
normalized volume.

\hypertarget{o10strength:proof-of-the-extension}{%
\subsubsection{Proof of the extension}\label{o10strength:proof-of-the-extension}}

Put \(P=\nabla v\), \(H=\nabla P\), \(q=|P|^2\), and use

\[
 w=\left(\frac{q+\delta}{R^2+q+\delta}\right)^\epsilon,
 \quad |dw|\sqrt q\le2\epsilon w|H|.
\]

The derivative-free Hessian interchange starts with

\[
 \int w|H|^2
 =\int w\operatorname{Re}\langle H_{ij},H_{ji}\rangle
      +\frac12\int w\sum_{i,j}|F_{ij}v|^2.
\]

Integrating once and commuting derivatives on the covector-valued \(P\)
produces the potential term \(\int w\langle P,\nabla h\rangle\), the
Ricci and connection-curvature pairings with \(P^2\), and the weight
derivative. Integrate this potential term once more:

\[
 \int w\operatorname{Re}\langle P,\nabla h\rangle
 =\int w|h|^2-\int\operatorname{Re}\langle dw\otimes P,h\rangle.
\tag{15}
\]

Boundedness of \(h\), the displayed weight derivative estimate, and
Young's inequality absorb the resulting \(C\epsilon\int w|H|\) term into
a small multiple of \(\int w|H|^2\) plus \(C\int w\). The weighted
bounded-section interpolation remains

\[
 \int wq^2\le C\int wq+C\int wq|H|
 \le C_1\left(\int w|H|^2+\int wq\right).
\tag{16}
\]

The curvature errors pair with \(wq\); Young gives a small multiple of
\(\int wq^2\) plus the normalized curvature-error squares. Thus the two
inequalities jointly give

\[
 \int w(|H|^2+q^2)\le C\int w(1+q)+o(V).
\tag{17}
\]

First take \(w=1\). Since \(\int q=\int\langle v,h\rangle\le CV\), this
proves the uniform fourth-power gradient and Hessian budget. Next use
the displayed slow weight: \(w\le R^{-2\epsilon}(q+\delta)^\epsilon\),
and the budget controls the average of \((q+\delta)^\epsilon(1+q)\) for
\(\epsilon\le1\). Let \(\delta\downarrow0\). This proves (14). The proof
uses bounded \(h\), not its derivative.

For the present proof, first use only the \textbf{unweighted} part of
this calculation. It works with a uniformly bounded normalized L2
curvature norm; the curvature-error term is bounded rather than small.
Equations (11)--(12) and (2a) therefore give

\[
 \|\nabla^2v\|_{2,\rm avg}+\|\nabla v\|_{4,\rm avg}\le C.
\tag{17a}
\]

This step does not assert uniform integrability of fourth-power
gradients.

\hypertarget{o10strength:amplitude-and-phase-control}{%
\subsubsection{6.1 Amplitude and phase
control}\label{o10strength:amplitude-and-phase-control}}

Put \(a=|v|^2\) and \(j_v=\operatorname{Im}\langle v,\nabla v\rangle\).
Since v and Tv are bounded,

\[
 \Delta a=2\operatorname{Re}\langle v,Tv\rangle-2|\nabla v|^2
\]

is bounded in normalized L2 by (17a). Pair it with \(a-m_*\to0\) in L2
to get

\[
 \|da\|_{2,\rm avg}^2
 =\fint(a-m_*)\Delta a\longrightarrow0.
\tag{17b}
\]

The strong H1 approximation and (2b) give
\(N_v=(\nabla v)^{1,0}_{J_+}\to0\) in L2. The elementary complex
decomposition of \(\langle v,\nabla v\rangle\) yields, away from rho
zero,

\[
 |j_v|\le C\bigl(|v|\,|N_v|+|da|\bigr).
\]

The low-rho sets have relative volume tending to zero. Their
gradient-square integral tends to zero by (17a), so they cause no loss.
Consequently

\[
 \|j_v\|_{2,\rm avg}\longrightarrow0.
\tag{17c}
\]

\hypertarget{o10strength:the-hopf-form-forces-an-exact-average-weyl-pairing}{%
\subsubsection{6.2 The Hopf form forces an exact average Weyl
pairing}\label{o10strength:the-hopf-form-forces-an-exact-average-weyl-pairing}}

Let \(Q_v=-iq_-(v)\), an ordinary real ASD two-form. The Hopf identities
in the manuscript's conventions are

\[
 |Q_v|^2=a^2/8,\qquad
 |\nabla Q_v|^2=\frac a2|\nabla v|^2-\frac12|j_v|^2,
 \qquad
 \delta Q_v=-\frac12j_v+E_v,
 \quad |E_v|\le C|v|\,|Dv|.
\tag{17d}
\]

Thus \(\delta Q_v\to0\) in L2, and ASD gives
\(\|dQ_v\|_2=\|\delta Q_v\|_2\to0\). From (13a),

\[
 \fint|\nabla v|^2=\fint\operatorname{Re}\langle v,Tv\rangle
 \longrightarrow2m_*.
\]

The factor a can be replaced by m* in its product with \(|\nabla v|^2\):
its L2 error tends to zero and the latter has a bounded L2 norm by
(17a). Hence

\[
 \fint|Q_v|^2\to m_*^2/8,\qquad
 \fint|\nabla Q_v|^2\to m_*^2.
\tag{17e}
\]

The ASD Bochner identity at s = −4 gives

\[
 2\fint\langle W^-Q_v,Q_v\rangle
 =\fint|\nabla Q_v|^2-\frac43\fint|Q_v|^2
 -\fint(|dQ_v|^2+|\delta Q_v|^2).
\]

Therefore

\[
 \boxed{\fint\langle W^-Q_v,Q_v\rangle\longrightarrow
                       \frac{5m_*^2}{12}.}
\tag{17f}
\]

For a trace-free symmetric operator on a three-dimensional Euclidean
space, \(\langle Wx,x\rangle\le\sqrt{2/3}|W|\,|x|^2\). Because a is
bounded and converges in measure to m*, the fourth moment of Q\_v tends
to \(m_*^4/64\). Applying Hölder to this sharp pointwise bound and using
(17f) shows

\[
 \liminf b_j\ge50/3.
\tag{17g}
\]

Together with (2a), this forces \(b_j\to50/3\).

\hypertarget{o10strength:equality-forces-a-smooth-bounded-full-weyl-model}{%
\subsubsection{6.3 Equality forces a smooth bounded full-Weyl
model}\label{o10strength:equality-forces-a-smooth-bounded-full-weyl-model}}

There is no need to divide by \(|Q_v|\) or choose eigenvectors. Put
\(q_0^2=m_*^2/8\) and define the globally smooth bounded trace-free
operator

\[
 W_0=\frac5{q_0^2}
       \left(Q_v\otimes Q_v-\frac{|Q_v|^2}{3}I_{\Lambda^-}\right).
\tag{17h}
\]

The exact algebraic identities and (17f) give

\[
 \fint|W_0|^2\longrightarrow50/3,
 \qquad\fint\langle W^-,W_0\rangle\longrightarrow50/3.
\]

Since \(b_j\to50/3\),

\[
 \boxed{\fint|W^--W_0|^2\longrightarrow0.}
\tag{17i}
\]

This proves the previously postponed bounded-model decomposition (4):
include W\_0, the scalar block, the actual bounded self-dual determinant
curvature and the bounded ASD model K in F\_0. The error consists of W−
minus W\_0, the trace-free Ricci block, and f minus K.

Now, and only now, apply the slow-weight conclusion (14) to v. It gives
uniform integrability of \(|\nabla v|^4\), exactly as required for the
Chern localization below.

\hypertarget{o10strength:the-chern-transgression-now-has-vanishing-normalized-cost}{%
\subsection{7. The Chern transgression now has vanishing normalized
cost}\label{o10strength:the-chern-transgression-now-has-vanishing-normalized-cost}}

Choose \(d>0\) with \(4d^2<m_*\). The standard rank-two transgression
applied to \(v/|v|\), with a smooth amplitude cutoff equal to one near
zero and zero above \(2d\), gives

\[
 |c_2(E)[M]|
 \le C_d\left[
 \int_{\{|v|<2d\}}|F^E|^2+
 \int_{\{d<|v|<2d\}}|\nabla v|^4\right].
\tag{18}
\]

The cutoff vanishes before any division by \(|v|\) is made at a zero, so
this is a global identity/estimate with all zero-set flux retained. It
does not require transverse zeros.

By (13), both sets in (18) have relative volume tending to zero. The
curvature squares are uniformly integrable by (4). The derivative fourth
powers are uniformly integrable by (14). Divide (18) by \(V\):

\[
 \boxed{c_2(E)[M]/V\longrightarrow0.}\tag{19}
\]

On the other hand, the actual Chern--Weil formula is

\[
 16\pi^2\frac{c_2(E)[M]}V
 =\fint\left[
 2|W^-|^2+\frac{s^2}{24}
 -\frac12|\operatorname{Ric}^{\circ}|^2
 +9|F_A^+|^2-9|F_A^-|^2\right].
\]

The actual endpoint limits are \(\fint|F_A^+|^2\to8/9\),
\(\fint|F_A^-|^2\to2/9\), \(s=-4\),
\(\fint|\operatorname{Ric}^{\circ}|^2\to0\), and the spectral argument
gives \(\fint|W^-|^2\to50/3\). Therefore

\[
 \boxed{\frac{c_2(E)[M]}V\longrightarrow
                   \frac{5}{2\pi^2}>0.}\tag{20}
\]

Equations (19) and (20) contradict each other. This proves the sequence
exclusion.

The auxiliary section \(v\) need not retain the actual cubic definition.
Its sole role in (18) is as a genuine section of the same bundle \(E\).
The actual field equations have already supplied its bounded approximate
eigenmode \(u\), while the resolvent supplies the critical regularity.
No field has been substituted into a monopole equation that it does not
satisfy.

\hypertarget{o10strength:the-universal-excess-gap-and-the-density-corollary}{%
\subsection{8. The universal excess gap and the density
corollary}\label{o10strength:the-universal-excess-gap-and-the-density-corollary}}

The excluded sequence was precisely (2a). Therefore the nonnegative left
side of (0) has a universal positive infimum. This proves the primary
excess-energy gap.

Chern--Gauss--Bonnet and the actual endpoint amplitude/Hodge bounds
give, under chi/A at most B,

\[
 \limsup\fint|W^-|^2\le\frac{16B}{9}-\frac23.
\tag{21}
\]

At \(B=39/4\), the right side equals \(50/3\), already excluded. There
is a nonzero absolute interval above this B as well. Indeed, if every
such interval admitted endpoint sequences, choose one with
\(B_n\le39/4+1/n\), then choose sufficiently deep terms to make beta
tend to infinity, kappa tend to zero, and the finite-error terms in (21)
tend to zero. This gives the forbidden sequence (2a). Equivalently, one
may use the positive margin in (0) directly.

For this fixed larger density interval, failure of the all-parameter
scalar-minus-four estimate would produce
\(\kappa_j\beta_j^2/(\beta_j-2)^2\to0\). The established
finite-parameter theorem, with the class-square cost removed, excludes
bounded beta. Hence beta tends to infinity and kappa tends to zero, a
contradiction. Finally the exact conformal comparison

\[
 \widehat\kappa\le\frac{\kappa}{1-\sqrt\kappa-\kappa/16}
\]

transfers the gap to every original admissible metric. The small-kappa
denominator is uniformly positive; metrics with kappa bounded below
satisfy (1) after reducing its constant. The actual supplier, class
square, positive budget, and Euler-density hypothesis are preserved by
the original comparison. This proves (1).

\hypertarget{o10strength:a-global-topological-cost-form-valid-on-every-original-admissible-metric}{%
\subsubsection{A global topological-cost form, valid on every original
admissible
metric}\label{o10strength:a-global-topological-cost-form-valid-on-every-original-admissible-metric}}

The density corollary can equivalently be packaged as another uniform
integral obstruction, without imposing a separate density restriction:

\[
 \boxed{\fint|\operatorname{Ric}+g|^2+\frac{D_\beta}V
 +\left(\frac{\chi(M)}A-\frac{39}{4}\right)_+
 \ge c'\frac{(\beta-2)^2}{\beta^2}\quad(\beta\ge2),}
\tag{22}
\]

where \(c'=\min\{c_*,\delta_*\}>0\). To verify it, put
\(f_\beta=(\beta-2)^2/\beta^2\le1\). If chi/A is at most
\(39/4+\delta_*\), use (1); otherwise the topology-excess term is larger
than \(\delta_*\ge\delta_* f_\beta\). Both other terms are nonnegative.
This proves (22) on the original admissible metric without needing a
conformal comparison for its opposite-Weyl energy.

Formula (22) preserves the meaning of the result: either the
Einstein/parameter defect or the explicit topological excess pays a
uniform amount. It does \textbf{not} delete that excess term from
unrestricted O10.

\hypertarget{o10strength:what-this-proves-and-what-remains}{%
\subsection{9. What this proves and what
remains}\label{o10strength:what-this-proves-and-what-remains}}

The theorem extends the fully controlled Euler-density range from the
earlier strict bound B \textless{} 3/8 to an absolute open interval
beyond 39/4. Its main statement (0) is stronger than a density-only
corollary: an actual weak zero endpoint must carry strictly more than
50/3 of normalized opposite-Weyl energy.

It does not prove an upper bound on that energy for arbitrary topology.
For sequences with a strictly larger opposite-Weyl energy, the
polynomial model (17h) need not have a vanishing L2 fitting error. The
Chern localization can then retain curvature and fourth-power derivative
concentration. Removing this excess-energy alternative remains the
unclosed part of general O10. The subsequent first-prolongation theorem
is a further valid input toward that problem, but is not part of the
proof above.
